\documentclass[12pt,a4paper]{article}
\usepackage[top=2.3cm,bottom=2.3cm,left=2.6cm,right=2.6cm]{geometry}
\usepackage[T1]{fontenc}
\usepackage{microtype}
\usepackage{parskip}
\usepackage{pdfpages}
\usepackage{xcolor}
\usepackage[colorlinks=true,urlcolor=blue,linkcolor=black]{hyperref}
\newcommand{\nrowslisted}{58}
\newcommand{\nrowsdropped}{2}
\newcommand{\mincutoff}{1}
 % \nrowslisted, \nrowsdropped, \mincutoff (written by changes_report.py)
\pagestyle{plain}
\setlength{\parskip}{0.55em}
\linespread{1.15}
\newenvironment{foldable}{}{}

\newcommand{\authorname}{Tri Nhan Vo}
\newcommand{\degree}{Doctor of Philosophy}
\newcommand{\submissiondate}{October, 2026}

\begin{document}

\begin{center}
  {\LARGE\bfseries Response to the Examiners' Reports \par}
  \vspace{0.5cm}
  {on the Thesis for the degree of \degree \par}
  \vspace{0.3cm}
  {\Large \bfseries Knowledge Distillation under Practical Constraints \\ for Models and Datasets \par}
  \vspace{0.3cm}
  {\large \authorname \par}
  {\submissiondate \par}
  \vspace{0.3cm}
\end{center}

\vspace{\baselineskip}

\begin{foldable}
  \textit{Dear Examiners,}
\end{foldable}

\begin{foldable}
  I sincerely thank you both for the time, dedication, and expertise you gave to examining my thesis.
  The reports were detailed and constructive, and they pointed to specific places where the thesis could say more precisely what it does and what its evidence supports.
  Revising my thesis in response has greatly improved its clarity, coherence, and overall quality.
\end{foldable}

\begin{foldable}
  \vspace{0.3cm}
  \textbf{\large Revision Summary}
\end{foldable}

\begin{foldable}
  The examination reports raised three main themes.
  The first theme is the \textbf{consistency and scope} of the claims.
  Several statements claimed somewhat more than the evidence fully supports.
  These included how my first body of work on \textit{DivBFKD} described the teacher's outputs, what high teacher confidence implies, how widely the state-of-the-art and ``first'' claims apply, and how closely a distilled corpus matches full-data performance.
  I checked each claim against the experiments and the chapters it appears in.
  I now state the access setting of each method exactly, name the baselines each result is compared with, and mark theoretical results as holding under stated assumptions.
  I also changed wording that depended on a proxy.
  One example is confidence used in place of closeness to the training distribution.
  The new wording now states closely what I measured.
\end{foldable}

\begin{foldable}
  The second theme is how accurately I described the \textbf{mechanisms}.
  The reports asked me to separate structural diversity from semantic diversity in \textit{DIPKD}.
  Their comments on the primer student and on TAKE pointed the same way: the primer student's outputs act as surrogate supervision, and \textit{TAKE}'s output is a selected surrogate corpus rather than a subset of the original data.
  I revised the method descriptions, the background chapter, and the concluding chapter so that they use the same terms and describe the same mechanism.
\end{foldable}

\begin{foldable}
  The third theme is the \textbf{practical constraints} that motivate the thesis.
  I first clarified what distributional fidelity means in each part of the thesis.
  I additionally introduced a new summary table in the concluding chapter, which compares the assumptions, information, mechanism, results, and limitations of the three methods.
  The table has put the research works in a comparable qualitative medium, and greately clarified the thesis' scope of work.
\end{foldable}

\begin{foldable}
  I would also like to explain, with respect, the places where I made a different choice.
  Distributional fidelity is not yet developed into a formal framework or a quantitative measure here, and this is restated as an open problem.
  Compute and teacher-query costs were not newly measured, so they are discussed qualitatively and marked as not fully measured.
  In TAKE, the embedding model used was fixed throughout, and thus reported as not ablated as a limitation.
  Also, QQP remains among the sentence-pair tasks, with a short note on its nature.
  Finally, claims of novelty and ``first'' are kept, since they held at the time of each publication.
  Your view on these choices would be welcome.
  Aside from these matters, I have also made quality-of-life changes, such as fixing minor typos and errors throughout to improve readability.
\end{foldable}

\begin{foldable}
  \vspace{0.3cm}
  \textbf{\large Tracked changes}
\end{foldable}

\begin{foldable}
  To make the changes easy to track, I also include (1) a \textit{Summary of Changes} table in the following pages of this response, and (2) an \textit{Editorial Version} of the thesis after that table, with in-line formatting of original and revised texts.
  They are generated from the two versions of the thesis and lists every segments that changed, appear in red strikethrough and inserted words in blue underline.
  Minor changes are skipped to avoid cluttering the table with insignificant edits.
\end{foldable}

\begin{foldable}
  To automatically and unbiasedly generate these ``\textit{diff-documents}'', I wrote a script that compares the difference and score them by a heuristic word-level \textit{Levenshtein distance} across a sentence or adjacent sentences.
  In short, inserting, deleting, or replacing a whole word costs one, a common function word (a/an/the, prepositions, etc.) costs 0.25.
  A segment's score is the sum over its words.
  The same rule applies to every sentences, and has been tested for reproducibility.
\end{foldable}

\begin{foldable}
  The Summary of Changes table, starting on the next page, lists every change that scores at least~\mincutoff, which means every change that alters at least one word.
  A higher cutoff would shorten the table, but it would hide some short but meaningful edits that addresses your concerns.
  That gives \nrowslisted{} segments with non-trivial changes.
  Please use the Editorial Version to read all changes in their original context, which follows after the table.
\end{foldable}

\begin{foldable}
  Thank you again for your careful reading and your guidance.
\end{foldable}

Sincerely,\\
Tri Nhan Vo

\clearpage
\includepdf[pages=-,fitpaper=true,pagecommand={\thispagestyle{empty}},picturecommand={\put(421,22){\makebox(0,0){\thepage}}}]{../../others/diff-workflow/output/summary-of-changes.pdf}
\includepdf[pages=-,fitpaper=true]{../../others/diff-workflow/output/main-diffexamined.pdf}

\end{document}
